\documentclass[12pt]{article}
\usepackage{amsmath,amssymb,amsthm,hyperref,booktabs,xcolor,enumitem}
\usepackage[margin=1in]{geometry}

\newtheorem{theorem}{Theorem}
\theoremstyle{definition}
\newtheorem{definition}[theorem]{Definition}
\newtheorem{remark}[theorem]{Remark}

\definecolor{high}{RGB}{180,50,50}
\definecolor{med}{RGB}{180,130,30}
\definecolor{low}{RGB}{80,130,80}

\title{\textbf{Open Conjecture Prioritization: Post-Verification Sprint}\\[0.4em]
\large After `analytic\_finite\_zeros\_compact` Proof + MUL Resolution}
\author{Arturo R.\ Almaguer}
\date{2026-04-21}

\begin{document}
\maketitle

\begin{abstract}
After completing the Lean verification sprint (21 theorems verified, 1 more proved this
session: \texttt{analytic\_finite\_zeros\_compact}), resolving CONJ\_MUL\_GEN\_TIGHT
($\mathrm{SB}(\mathrm{mul}) = 3$), and identifying the exact sorry structure in T01, we
re-rank open conjectures by mathematical leverage and Lean-feasibility.
The new top priority is \textbf{eml\_tree\_analytic}: a 2-hour Lean engineering task
that would close 2 more theorems (T01 analyticity + T48 depth barrier) and enable the
full EDB construction.
\end{abstract}

%──────────────────────────────────────────────────────────────────────────────
\section{Current Verified Count and Sorry Structure}

\subsection{Fully Lean-Verified (0 sorry): 17 theorems}
\begin{enumerate}[noitemsep]
  \item SB(neg/add/sub/mul/div) $\geq 2$ — 5 lower bounds
  \item exp/mul/pow/recip/sqrt node counts — 5 upper bounds
  \item v5.1 overcounting correction (16n → 15n)
  \item Euler Gateway, Euler Identity, EML-0 $\subsetneq$ EML-1
  \item sin has infinitely many zeros
  \item \textbf{[NEW]} analytic\_finite\_zeros\_compact (proved this session)
  \item TheoremRegistry rfl check
\end{enumerate}

\subsection{Sorry'd (4 genuine open steps)}

\begin{center}
\begin{tabular}{llll}
\toprule
ID & Theorem & Sorry type & Effort \\
\midrule
C1 & \texttt{eml\_tree\_analytic} & Lean engineering & 2h \\
C2 & \texttt{sin\_not\_in\_eml} & O-minimal theory & months \\
C3 & \texttt{depth1\_finite\_zeros\_real} & depends on C1 & 30min after C1 \\
C4 & \texttt{sin\_not\_in\_real\_EML\_k} & depends on C2 & months \\
\bottomrule
\end{tabular}
\end{center}

Note: C3 is unlocked by C1 (just applies eml\_tree\_analytic + analytic\_finite\_zeros\_compact).
C4 is unlocked by C2 (genuine open).

%──────────────────────────────────────────────────────────────────────────────
\section{Ranked Open Items}

\subsection{Priority 1 — \texttt{eml\_tree\_analytic} [2h Lean, no math uncertainty]}
\label{sec:p1}

\textbf{Statement}: $\forall t : \texttt{EMLTree},\ \texttt{AnalyticOnNhd}\ \mathbb{R}\ t.\texttt{evalReal}\ (0, \infty)$.

\textbf{Why top}: Pure formalization — the math is known. Unblocks C3, tightens T01.

\textbf{Proof sketch}:
\begin{verbatim}
induction t with
| const c => exact analyticOnNhd_const.mono ...
| var     => exact analyticOnNhd_id.mono ...
| ceml t1 t2 ih1 ih2 =>
    -- evalReal x = exp(t1.evalReal x) - (Complex.log t2.eval ↑x).re
    apply AnalyticOnNhd.sub
    · -- exp ∘ t1.evalReal: analytic by ih1 + analyticOnNhd_exp.comp
    · -- (Complex.log t2.eval ↑x).re: analytic on (0,∞) where t2.eval positive
      -- Key: need to show t2.eval ↑x ≠ 0 for x ∈ (0,∞) or handle via re of Complex.log
\end{verbatim}

\textbf{Main import needed}: \texttt{Mathlib.Analysis.Analytic.Composition},
\texttt{Mathlib.Analysis.SpecialFunctions.Complex.Analytic} (or similar for Complex.log analyticity).

\textbf{Impact}: Closes 2 sorries (C1, C3). Lifts sorry count from 4 to 2.

---

\subsection{Priority 2 — CONJ\_DIV\_GEN\_TIGHT: SB(div) = 3}

\textbf{Method}: Clone \texttt{mul\_gen\_tight\_2node\_search.py}, replace target with $x/y$, run.

\textbf{Effort}: 30 minutes.  \textbf{Expected result}: SB(div, general) = 3.

\textbf{Impact}: Closes another $[2,6]$ gap. Python-certified.

---

\subsection{Priority 3 — CONJ\_BOUNDARY\_DECIDABLE}

\textbf{Statement}: SC + EDB $\Rightarrow$ ELC-membership decidable.

\textbf{Status}: The EDB positive-domain part ($B^+(n) = n$) is conceptually proved.
The EDB general-domain part requires $Z(n)$ (quantitative zero bound for EML trees on $\mathbb{R}$),
which is the same obstruction as C2 (o-minimal structure).

\textbf{Impact}: Landmark theorem. Requires C2 first.

---

\subsection{Priority 4 — CONJ\_TRIG\_DEPTH\_TOWER L1}

\textbf{Statement}: $\sin(1) \notin \varepsilon(\mathbb{R})$ (unconditional L1).

\textbf{Status}: Requires a new algebraic tool. The AIL (Algebraically Isolated Limit)
argument gives $\tan(1) \notin \varepsilon(\mathbb{R})$ (proved), but $\sin(1)$ needs a
different approach since sin is not a ratio of exp values.

\textbf{Effort}: Open mathematical problem. Likely requires Schanuel-level machinery.

---

\subsection{Priority 5 — GAIL}

\textbf{Statement}: $\alpha \notin \varepsilon(\mathbb{R}) \Rightarrow \sin(\alpha) \notin \varepsilon(\mathbb{R})$.

\textbf{Status}: Unconditionally open. Follows from Schanuel's Conjecture.
Proving GAIL unconditionally is a major open problem in transcendence theory.

---

\subsection{Priority 6 — eml\_tree\_analytic (General Domain)}

\textbf{Statement}: Every EML tree is analytic on all of $\mathbb{R}$ (or on $\mathbb{R} \setminus Z_t$
for a computable isolated set $Z_t$).

\textbf{Status}: This is the fully general form of C1. The issue is that
$\texttt{ceml}(t_1, t_2)$ can be non-analytic where $t_2.\texttt{eval}\ x = 0$.
This requires tracking the zero set of each sub-tree, which is an inductive computation.

\textbf{Impact}: Would give the quantitative zero bound B(k) needed for C2, closing the
circle to T01 Part D.

%──────────────────────────────────────────────────────────────────────────────
\section{Concrete Attack Plan: Priority 1 (eml\_tree\_analytic)}

\subsection{Step 1: Required Imports}
\begin{verbatim}
import Mathlib.Analysis.Analytic.Composition
import Mathlib.Analysis.SpecialFunctions.ExpDeriv
-- Check if available:
import Mathlib.Analysis.SpecialFunctions.Complex.AnalyticLog
\end{verbatim}

\subsection{Step 2: Helper Lemmas}

\textbf{Lemma 1}: $\texttt{Complex.exp}$ is analytic on $\mathbb{R}$ (as a real function).
\begin{verbatim}
lemma real_exp_analyticOnNhd :
    AnalyticOnNhd ℝ Real.exp Set.univ := ...
-- Use: Real.analyticAt_exp : AnalyticAt ℝ Real.exp x
-- Then: AnalyticOnNhd ℝ Real.exp univ from this
\end{verbatim}

\textbf{Lemma 2}: $x \mapsto (t_1.\texttt{eval}\ \uparrow x).\texttt{re}$ is analytic on $(0,\infty)$.
This is exactly \texttt{eml\_tree\_analytic t1} by IH.

\textbf{Lemma 3}: $x \mapsto (Complex.\texttt{log}\ (\uparrow x)).\texttt{re}$ is analytic on $(0,\infty)$.
This equals $\mathbb{R}.\texttt{log}\ x$ for $x > 0$, which is analytic on $(0,\infty)$.
\begin{verbatim}
lemma real_log_analyticOnNhd_pos :
    AnalyticOnNhd ℝ Real.log (Set.Ioi 0) :=
  Real.analyticOnNhd_log.mono ...  -- or analogous
\end{verbatim}

\subsection{Step 3: Inductive Step for ceml}

\begin{verbatim}
| ceml t1 t2 ih1 ih2 =>
  simp only [EMLTree.evalReal, EMLTree.eval]
  apply AnalyticOnNhd.sub
  · -- exp(t1.evalReal x): compose ih1 with Real.exp analytic
    exact real_exp_analyticOnNhd.comp ih1 ...
  · -- (Complex.log (t2.eval ↑x)).re on (0,∞):
    -- For t2 = const c: constant, analytic
    -- For t2 = var: Complex.log ↑x .re = Real.log x, analytic on (0,∞)
    -- For t2 = ceml ...: ih2 gives AnalyticOnNhd t2.evalReal (Ioi 0)
    --   and t2.evalReal x > 0 for x ∈ (0,∞) would give analytic log
    --   BUT we don't know t2.evalReal > 0 in general
    -- Workaround: Complex.log is analytic on ℂ \ ℝ≤0; for t2.eval ↑x (real x > 0),
    --   if t2.evalReal x > 0 then Complex.log ↑(t2.evalReal x) = Real.log (t2.evalReal x)
    --   If t2.evalReal x ≤ 0: Complex.log puts branch cut, real part = Real.log |t2.evalReal x|
    --   The real part function is analytic EVERYWHERE on ℝ (discontinuity only at x=0 in
    --   the imaginary direction).
    --
    -- ACTUAL RESOLUTION: (Complex.log ↑(t2.evalReal x)).re = Real.log |t2.evalReal x|
    --   wherever t2.evalReal x ≠ 0, and = 0 at t2.evalReal x = 0.
    --   This equals Real.log ∘ |·| ∘ t2.evalReal, where:
    --   - |·| is analytic away from 0
    --   - Real.log is analytic on (0,∞)
    --   So the composition is analytic wherever t2.evalReal ≠ 0.
    --   On (0,∞), if t2.evalReal has isolated zeros, the term is analytic almost everywhere.
    --   The cleanest fix: assume t2.evalReal x ≠ 0 for x ∈ domain — track as hypothesis.
    sorry
\end{verbatim}

\subsection{Difficulty Assessment}

The core difficulty in step 3 is the log term. For the specific use case
(trees applied to positive real inputs), the log argument is typically positive,
making the proof feasible. But for full generality, domain tracking per node is needed.

\textbf{Estimate}: 2–3 hours of Lean engineering if one focuses on the case
$t_2.\texttt{evalReal}\ x > 0$ for $x > 0$ (sufficient for the sin barrier application).

%──────────────────────────────────────────────────────────────────────────────
\section{Summary}

\begin{center}
\begin{tabular}{lllll}
\toprule
Priority & ID & Statement & Effort & Next step \\
\midrule
1 & eml\_tree\_analytic & AnalyticOnNhd on (0,∞) & 2–3h Lean & write proof \\
2 & DIV-TIGHT & SB(div) = 3 & 30min script & run script \\
3 & DECIDABLE & SC+EDB → decidable & months & needs C1+C2 \\
4 & TRIG-L1 & sin(1) ∉ ε(ℝ) & open & new tool \\
5 & GAIL & sin preserves non-ELC & open & SC needed \\
\bottomrule
\end{tabular}
\end{center}

\textbf{Recommended immediate next step}: Prove \texttt{eml\_tree\_analytic} for
the constant and var base cases first (trivial), then handle the ceml step assuming
$t_2.\texttt{evalReal}\ x > 0$ for $x > 0$ (the main use case).
This would reduce the sorry count from 4 to 2 and lift the verified theorem count
from 17 to 19.

\end{document}
